\section{Discussion}
\label{sec:discussion}

\paragraph{The error is in the source model, and both adaptation methods inherit it.} The gap between seen and unseen
bearings is present in the source model before any adaptation step (median 35.1 points) and is of the same size afterwards.
SDALR's self-training leaves it essentially unchanged; SHOT narrows it, but by giving up accuracy on the seen bearings rather
than by gaining on the unseen ones; a forest that never adapts shows it too. A source-free method can only refine what its
source model already separates, and on an unseen specimen the source model separates the wrong thing. This outcome was not
foregone. Self-training could have repaired the gap: the oracle arm shows that the same network, adapted with correct
pseudo-labels, recovers a median 82\,\% of it, so the network can be steered to most of the missing accuracy. The two methods' own
pseudo-labels simply carry the source model's errors forward. The question an
SFDA result on a bearing rig has to answer is therefore not ``how much does adaptation gain?'' but ``what did the source model
learn that transfers to a specimen it has not seen?''.

\paragraph{Under-sampling, not a fixed identity shortcut.} Individual specimens are identifiable from ten time statistics,
but that alone explains little: healthy specimens are identifiable too and transfer almost perfectly, and adding source
specimens raises a forest's unseen-bearing accuracy from 46 to 75\,\% with the test set fixed. Faulty specimens differ from each other in
what matters for the label --- the size, shape and mechanism of their damage (Table~\ref{tab:bearings}) --- and three of them per
class do not span that variation. This reading has a practical consequence the identity reading lacks: the remedy is more
diverse source specimens, as Vieira et al.\ found for supervised models \cite{vieira2026}, rather than an adversary that hides
specimen identity. A class-conditional identity adversary may still help at the margin, but the curve says diversity is where
most of the missing accuracy lies. The two readings are not exclusive: with one or two specimens per class, features
particular to those specimens separate the classes as well as any shared fault feature, and only a more diverse source makes
the shared features the cheaper solution. In that sense source diversity breaks the reliance on specimen-specific features;
the curve shows how much accuracy that is worth. The curve was measured with a forest; whether SDALR or SHOT follow it is
not measured.

\paragraph{Physics gating is bounded by detectability.} The safe gates help, consistently and modestly: a median 4.5 points
over the ten splits, essentially never costing anything. Three limits follow from the measurements. First, a filter can only
act through correct labels on the recordings it certifies, and the gain lies almost entirely there; the gate does not teach
the model anything that carries over to recordings it is silent on. Second, the envelope gates certify only the specimens with
the largest damage, and every envelope-based gate we tried certifies the same ones; detectability is a property of the damage, not of
the threshold. Third, an oracle filter shows that correct pseudo-labels on every recording would recover most of the gap, so
the room is there --- but it is room that envelope evidence cannot fill for small, early-stage damage, which is the damage a
monitoring system most needs to detect. A kinematic representation leaks less for the same reason: it relies only on what is
visible, and on this rig that is the large damage.

\paragraph{Safety is a property to measure, not to assume.} The raw-spectrum band rule accepts \SI{45}{\percent} of healthy
recordings as faulty; the envelope rules accept at most \SI{2}{\percent}. Inside the loop the safe gates never cost more than
1.5 points on a split, whereas the unsafe rule moved accuracy by $+5.4$ or $-13.2$ points depending on which bearings were in
the target. A module that injects confident wrong labels can look helpful on one target and harmful on another, and end
accuracy on one leaky protocol cannot tell a safe prior from an unsafe one. False acceptance on healthy recordings can.

\paragraph{A shallow baseline belongs in every such comparison.} On bearing-disjoint targets a forest on ten time-domain
statistics --- no adaptation, no target data, no deep network --- is more accurate than adapted SDALR in 8 of 10 splits. It is
not a better diagnostic method; it collapses too. It is a calibration point that any claim of transferable fault structure
should clear.

\paragraph{Is this leakage, or a different task?} Adapting to the same physical bearings at a new operating condition is a
legitimate task for a permanently instrumented asset. Two scenarios are conflated in the current literature. In the first the
monitored bearing is present in training, and same-bearing evaluation estimates performance honestly. In the second ---
a replacement bearing, a sister machine, a fleet --- the specimen at test time was never seen, and same-bearing evaluation
over-estimates performance by 22--36 points (median over splits, across the methods measured here). The problem is not the
published protocol in itself but reading its number as evidence for the second scenario. We use ``leakage'' in the sense of
Kapoor and Narayanan \cite{kapoor2023leakage}: information shared between training and evaluation that inflates the estimate
relative to the deployment the claim addresses.

\paragraph{Threats to validity.} The leaky and clean targets differ in their bearing sets; the within-bearing contrast holds
the specimen fixed and gives the same answer. The splits share bearings from a pool of 17, so their $p$-values are conditional
on this pool. The two folds are asymmetric (fold~B has two outer-race source bearings) and ran in two jobs with different
source checkpoints, which is why the folds are descriptive. The HUST replication changes bearing size along with the
specimen. Finally, the windows given to the networks cannot carry the kinematic signature; the kinematic-feature forest shows
what changes when the input can.
